% ============================================================================
\section{Formulación de la hipótesis}

\subsection{Hipótesis general}

La incorporación de una representación temporal generada por un modelo QRC mejora la predicción fuera de muestra de las revisiones del PBI peruano respecto del mejor benchmark clásico preespecificado, siempre que la estructura de las revisiones contenga dependencias no lineales explotables y la muestra permita estimarlas sin sobreajuste.

\subsection{Hipótesis específicas}

\begin{hypothesis}[Estructura predecible de las revisiones]
	Las revisiones del PBI presentan al menos una propiedad sistemática —sesgo medio, persistencia o asociación con información disponible en la primera publicación— que puede ser identificada mediante pruebas descriptivas e inferenciales apropiadas.
\end{hypothesis}

\begin{hypothesis}[Desempeño incremental de QRC]
	En la evaluación temporal fuera de muestra, QRC reduce la pérdida predictiva frente al mejor benchmark clásico preespecificado y frente al reservorio clásico de plausibilidad, con una mejora cuya magnitud e incertidumbre se reportan explícitamente.
\end{hypothesis}

\begin{hypothesis}[Dependencia respecto de las condiciones del problema]
	El desempeño relativo de QRC depende de características del proceso generador, en particular del tamaño muestral, la persistencia, la dimensionalidad efectiva y la razón señal-ruido.
\end{hypothesis}

\begin{hypothesis}[Relevancia económica condicionada]
	Si QRC produce una mejora predictiva estable y replicable, la información adicional puede reducir la pérdida de una regla de decisión económica bajo supuestos explícitos. Esta hipótesis es secundaria y no se considerará contrastada únicamente por una reducción del RMSE.
\end{hypothesis}

\subsection{Expresión funcional de las variables}

La variable dependiente principal fue la revisión a horizonte fijo:
\[
Y_t^{(H)}=y_t^{(v_H(t))}-y_t^{(v_0(t))}.
\]
Las variables explicativas fueron exclusivamente aquellas estrictamente observables en el origen \(v_0(t)\): la primera estimación publicada, las revisiones completadas de períodos previos en o antes de \(v_0(t)\), las tasas iniciales y revisiones rezagadas de las ramas productivas continuas, y las variables calendario documentadas. La especificación del vector de predictores quedó cerrada tras auditar la disponibilidad histórica de la base y antes de estimar los modelos predictivos.

\begin{table}[htbp]
	\centering
	\caption{Definición operativa de variables}
	\label{tab:variables}
	\footnotesize
	\renewcommand{\arraystretch}{1.15}
	\begin{tabular}{@{}>{\raggedright\arraybackslash}p{3.1cm}>{\raggedright\arraybackslash}p{2.3cm}>{\raggedright\arraybackslash}p{6.1cm}>{\raggedright\arraybackslash}p{2.5cm}@{}}
		\toprule
		Variable & Tipo & Definición operacional & Fuente auditada\\
		\midrule
		Revisión a horizonte \(H\) & Dependiente & Diferencia entre la publicación a horizonte fijo y la primera publicación del mismo período & Vintages auditados del PBI\\
		Primera estimación & Explicativa & Primera tasa de crecimiento publicada para el período & Archivo vintage BCRP\\
		Revisiones históricas & Explicativa & Revisiones de períodos anteriores completadas y observables en el origen & Archivo vintage auditado\\
		Indicadores sectoriales & Explicativa & Primeras publicaciones y rezagos de ramas productivas disponibles en cada origen & Archivo vintage sectorial\\
		Calendario de publicación & Control & Mes, trimestre y distancia temporal desde la primera publicación & Metadatos de releases\\
		\bottomrule
	\end{tabular}
	\fuente{elaboración propia.}
\end{table}
